\documentclass{amsart}
\usepackage[T1]{fontenc}

\addtolength{\textwidth}{2cm}
\addtolength{\hoffset}{-1cm}
\addtolength{\textheight}{2cm}
\addtolength{\voffset}{-1cm}

\renewcommand{\baselinestretch}{1.5}

\usepackage{amsmath,amsthm,hyperref,amssymb,mathtools,enumitem,tikz,tikz-cd}
\usepackage{xcolor}
% \usepackage{libertine,libertinust1math}
\usepackage{mathrsfs}
\usetikzlibrary{calc}
\usepackage{graphicx} % Required for inserting images
\numberwithin{equation}{section}
\makeatletter
\def\l@subsection{\@tocline{2}{0pt}{2.5pc}{5pc}{}}
\makeatother

\usetikzlibrary{arrows.meta,decorations.pathmorphing}

\title{Stokes theorem revisited}
\author{Oscar Harr and Florian Riedel}
\date{\today}

\newtheorem{thm}{Theorem}[section]
\newtheorem{mainthm}{Theorem}
\renewcommand{\themainthm}{\Alph{mainthm}}
\newtheorem*{thm*}{Theorem}
\newtheorem{lem}[thm]{Lemma}
\newtheorem{prop}[thm]{Proposition}
\newtheorem{cor}[thm]{Corollary}

\theoremstyle{definition}
\newtheorem{defn}[thm]{Definition}
\newtheorem{cons}[thm]{Construction}
\newtheorem{convent}[thm]{Convention}
\newtheorem{notn}[thm]{Notation}
\newtheorem{rmk}[thm]{Remark}
\newtheorem{fact}[thm]{Fact}
\newtheorem{exmp}[thm]{Example}

\DeclareMathOperator{\colimit}{\operatorname{colim}}
\DeclareMathOperator{\limit}{\operatorname{lim}}
\tikzset{%
    symbol/.style={%
        draw=none,
        every to/.append style={%
            edge node={node [sloped, allow upside down, auto=false]{$#1$}}}
    }
}

\newcommand{\Z}{\mathbb Z}
\newcommand{\Q}{\mathbb Q}
\newcommand{\R}{\mathbb R}

\newcommand{\PresMon}{\mathrm{Pr}_{\mathrm{st}}^{L,\otimes}}
\newcommand{\LinCatMon}[1]{\mathrm{LinCat}_{\mathrm{st},#1}^{\otimes_R}}
\newcommand{\LinCat}[1]{\mathrm{LinCat}_{\mathrm{st},#1}}
\newcommand{\Pres}{\mathrm{Pr}_{\mathrm{st}}^{L}}
\newcommand{\NcMot}{\mathrm{NcMot}}
\newcommand{\colimCats}{\mathrm{Pr}^L}
\newcommand{\unstLinCat}[1]{\mathrm{LinCat}_{#1}}
\newcommand{\Sp}{\mathrm{Sp}}
\newcommand{\Fin}{\mathrm{Fin}}
\newcommand{\lchaus}{\mathrm{LCHaus}}
\newcommand{\Ab}{\mathrm{Ab}}
\newcommand{\frBord}[2]{\mathrm{Bord}_{#1}^{\mathrm{fr}}(\R^{#2})}
\newcommand{\modules}[1]{\operatorname{Mod}_{#1}}

\newcommand{\support}{\operatorname{supp}}
\newcommand{\interior}{\operatorname{int}}

\newcommand{\Corr}{\operatorname{Corr}}
\newcommand{\Fun}{\operatorname{Fun}}
\newcommand{\Hom}{\operatorname{Hom}}
\newcommand{\iHom}{\underline{\operatorname{Hom}}}
\newcommand{\Aut}{\operatorname{Aut}}
\newcommand{\Map}{\operatorname{Map}}
\newcommand{\intHom}{\underline{\operatorname{Hom}}}

\newcommand{\term}{\mathrm{pt}}

\newcommand{\trace}{\operatorname{tr}}
\newcommand{\rig}[1]{#1^{\mathrm{rig}}}

\newcommand{\Shv}{\operatorname{Shv}}
\newcommand{\ccoShv}{\widehat{\operatorname{coShv}}}
\newcommand{\CAlg}{\operatorname{CAlg}}

\newcommand{\note}[1]{{\bf TODO: #1}}

\begin{document}
\begin{thm}
  Let $X$ be a compact Hausdorff space and let $f\colon X\to X$ be a continuous map. There is an equivalence
  $$
  \trace\left(
    f_?\curvearrowright\ccoShv (X;\Sp )
  \right)
  \simeq\Sigma^\infty_+|X^f|,
  $$
  natural in the pair $(X,f)$. 
\end{thm}
\begin{proof}
  Let $\frBord 1 1$ be the $\infty$-category of framed tangles in $\lbrack 0,1\rbrack\times\R$. By the one-dimensional tangle hypothesis, this is the free monoidal stable $\infty$-category on a single dualizable object. Hence
  $$
  \mathcal E = \Fun (\frBord 1 1^{\text{op}},\Sp )
  $$
  with the monoidal structure of Day convolution is the free presentably monoidal category on a single dualizable object. An $\mathcal E$-module in $\Pres$ is the same as a presentable stable $\infty$-category $\mathcal C$ with a distinguished strongly continuous endofunctor $T\colon\mathcal C\to\mathcal C$. Explicitly, this endofunctor is given by scalar multiplication with the object $\Sigma^\infty_+y(\term )\in\Fun (\frBord 1 1^{\text{op}},\Sp )$, where $y\colon\frBord 1 1^{\text{op}} \to\mathcal S$ is the Yoneda embedding. It follows from the rigidity of $\frBord 1 1$ that $\mathcal E$ is rigid. By a theorem of Arinkin--Gaitsgory--Kazhdan--Raskin--Rozenblyum--Varshavsky, a dualizable $\mathcal E$-module is the same as an $\mathcal E$-module whose underlying presentable stable category is dualizable. We are interested in the finitary localizing invariant
  $$
  \tau\colon\modules{\mathcal E}(\Pres )\to\Sp
  $$
  given by sending $\mathcal C$ to $\trace\left( \term\curvearrowright \mathcal C\right)$. Specifically, we want to determine the value of $\tau$ on $\ccoShv (X;\Sp )$. Here we can identify $\ccoShv (X;\Sp )$ with the $\mathcal E$-linearly enriched Hom into the $(\mathcal E,\mathcal E)$-bimodule $\Sp$. It then follows by a theorem of Efimov that the universal finitary $\mathcal E$-linear localizing invariant $\mathcal U_{\text{loc},\mathcal E}$ that there is an equivalence
  $$
  \mathcal U_{\text{loc},\mathcal E}\ccoShv(X;\Sp )
  \simeq\intHom_{\NcMot_{\mathcal E}}(\mathcal U_{\text{loc},\mathcal E}\Shv (X;\Sp ),\mathcal U_{\text{loc},\mathcal E}\Sp ),
  $$
  natural in $X$. \note{Finish this argument! There are still some details missing. It would be particularly clean if $\NcMot_{\mathcal E}$ has some description in terms of $\NcMot$, e.g. as $\NcMot^{\Delta^1}$.}
\end{proof}
\begin{prop}
  Let $X$ be a compact Hausdorff space and let $f\colon X\to X$ be a continuous map. Assume that $X^f$ is a finite collection of points. The diagram
  $$
  \begin{tikzcd}
    \Sp\arrow[r,"w^{\rm loc}_X"]\arrow[d,equal]
    &\ccoShv (X;\Sp )\arrow[r,"a^*"]\arrow[d,"f_?"]\arrow[dl,shorten >=4ex,shorten <=4ex,Rightarrow]
    & \Sp^{|X|}\arrow[d,"f_!"] \\
    \Sp\arrow[r,swap,"w^{\rm loc}_X"]
    &\ccoShv (X;\Sp )\arrow[r,swap,"a^*"]
    & \Sp^{|X|}
  \end{tikzcd}
  $$
  induces a commutative triangle $\mathbb S\to\Sigma^\infty_+ |X^f|\to\Sigma^\infty_+ L^f|X|$ \note{finish this}
\end{prop}
\note{Note that in the above }
\end{document}
